\section{Justificación}

La formación práctica en ciberseguridad es una necesidad relevante para estudiantes de Ingeniería Informática e Ingeniería de Sistemas, debido a que las amenazas digitales actuales exigen análisis, toma de decisiones y respuesta oportuna ante incidentes. En la asignatura de Seguridad de Sistemas, el enfoque principalmente conceptual y la cantidad elevada de estudiantes dificultan la práctica guiada y el seguimiento individual.

Si esta situación no se atiende, los estudiantes pueden finalizar su formación con dificultades para identificar riesgos digitales, priorizar medidas de protección y comprender las consecuencias de una respuesta inadecuada frente a amenazas como \textit{phishing}, malware o ingeniería social.

Por ello, el proyecto propone desarrollar un videojuego de simulación gamificada como herramienta tecnológica de apoyo al aprendizaje práctico. Mediante escenarios interactivos, retos progresivos, toma de decisiones y reportes de retroalimentación, el estudiante podrá practicar en un entorno controlado, seguro y accesible.

Los principales beneficiarios serán los estudiantes que cursan contenidos relacionados con seguridad informática, quienes contarán con un recurso complementario para aplicar conceptos y decisiones de defensa. También los docentes podrán utilizar la herramienta como apoyo para dinamizar la enseñanza y reforzar contenidos mediante actividades interactivas.

El proyecto es viable porque se plantea como un prototipo monojugador de alcance acotado, centrado en amenazas digitales fundamentales y mecánicas de simulación educativa, utilizando herramientas tecnológicas accesibles dentro del tiempo disponible para Taller de Grado I.
